% =============================
% ภาคผนวก
% =============================
% หมายเหตุ: ก่อนเริ่มภาคผนวก ระบบจะสร้างหน้าคั่นให้อัตโนมัติ
% ใช้คำสั่ง \appendiceschapter{ชื่อภาคผนวก} สำหรับแต่ละภาคผนวก
% ระบบจะกำหนดตัวอักษร ก, ข, ค ... และเลขหัวข้อ (ก.1, ก.1.1) ให้โดยอัตโนมัติ
%
% ภาพหน้าจอทั้งหมดคัดจากคู่มือการใช้งาน (manual/images/screens) เฉพาะหน้าจอของส่วนเว็บ
% หน้าจอของแชตบอตและระบบวางแผนการเดินทางไม่ได้นำมาใส่ในภาคผนวกนี้

% ภาพหน้าจอ: \appscreen[ความกว้าง]{ชื่อไฟล์ไม่รวมนามสกุล}{คำอธิบายภาพ}
\newcommand{\appscreen}[3][0.85]{%
\begin{figure}[H]
\centering
\includegraphics[width=#1\textwidth,keepaspectratio]{images/screens/#2.jpg}
\caption{#3}
\label{fig:app_#2}
\end{figure}}

% =============================
% ภาคผนวก ก
% =============================
\appendiceschapter{ภาพหน้าจอระบบที่พัฒนา}

ภาคผนวกนี้แสดงหน้าจอของระบบที่พัฒนาเสร็จแล้วในส่วนเว็บ เพื่อให้เทียบกับแบบร่างหน้าจอ (Wireframe) ในบทที่ 4 ได้ แบ่งเป็นหน้าจอของผู้ใช้งานทั่วไปและสมาชิก และหน้าจอของผู้ดูแลระบบ ข้อมูลที่ปรากฏในภาพ เช่น ชื่อ อีเมล และเบอร์โทรศัพท์ของผู้ใช้ เป็นข้อมูลตัวอย่างที่สร้างขึ้นสำหรับการจัดทำเอกสาร มิใช่ข้อมูลของผู้ใช้จริง

\section{หน้าจอของผู้ใช้งานทั่วไปและสมาชิก}

\subsection{หน้าแรก การสมัครสมาชิก และการเข้าสู่ระบบ}

ภาพที่ \ref{fig:app_guest-home} แสดงหน้าแรกที่รวมกิจกรรมแนะนำและสถานที่แนะนำไว้ในหน้าเดียว พร้อมแถบเมนูสำหรับผู้ใช้ที่ยังไม่เข้าสู่ระบบ ภาพที่ \ref{fig:app_guest-signup-form} ถึง \ref{fig:app_forgot-password} แสดงขั้นตอนสมัครสมาชิก เข้าสู่ระบบ และตั้งรหัสผ่านใหม่ (UC-15 และ UC-16)

\appscreen{guest-home}{หน้าแรกและแถบเมนูผู้ใช้งานทั่วไป}
\appscreen{guest-signup-form}{หน้าฟอร์มสมัครสมาชิก}
\appscreen{guest-login-form}{หน้าฟอร์มเข้าสู่ระบบ}
\appscreen{forgot-password}{หน้าลืมรหัสผ่าน}

\subsection{สถานที่ท่องเที่ยวและกิจกรรม}

ภาพที่ \ref{fig:app_places-list} และ \ref{fig:app_place-detail} แสดงรายการสถานที่ที่มีช่องค้นหาและตัวกรองหมวดหมู่ และหน้ารายละเอียดสถานที่ (UC-02 และ UC-04) ส่วนภาพที่ \ref{fig:app_events-list} และ \ref{fig:app_event-detail} แสดงรายการและรายละเอียดกิจกรรมที่มีโครงสร้างสอดคล้องกัน (UC-05)

\appscreen{places-list}{หน้ารายการสถานที่ท่องเที่ยว}
\appscreen{place-detail}{หน้ารายละเอียดสถานที่ท่องเที่ยว}
\appscreen{events-list}{หน้ารายการกิจกรรมและเทศกาล}
\appscreen{event-detail}{หน้ารายละเอียดกิจกรรม}

\subsection{พอยท์สะสมและการสแกน QR}

ภาพที่ \ref{fig:app_points-guest} แสดงหน้าพอยท์สะสมสำหรับผู้ที่ยังไม่เข้าสู่ระบบ ซึ่งบอกวิธีสะสมและแสดงสถานที่ที่มี QR ให้สะสมพอยท์ได้ ส่วนภาพที่ \ref{fig:app_points-member} และ \ref{fig:app_qr-scan} แสดงหน้าพอยท์ของสมาชิกและกล่องสแกน QR ที่ขอใช้กล้องและตำแหน่งของผู้ใช้ (UC-03 และ UC-08)

\appscreen{points-guest}{หน้าพอยท์สะสม กรณีผู้ใช้งานทั่วไป}
\appscreen{points-member}{หน้าพอยท์สะสมของสมาชิก}
\appscreen{qr-scan}{กล่องสแกน QR Code เพื่อสะสมพอยท์}

\subsection{บัญชีของสมาชิก}

หน้าบัญชีของสมาชิกแบ่งเป็น 5 แท็บ ได้แก่ ภาพรวม ข้อมูล ประวัติการสแกนและการแลก กิจกรรมของฉัน (คำขอเพิ่มกิจกรรม) และความปลอดภัย แสดงดังภาพที่ \ref{fig:app_profile-overview} ถึง \ref{fig:app_profile-security} (UC-17 และ UC-19)

\appscreen{profile-overview}{หน้าโปรไฟล์ แท็บภาพรวม}
\appscreen{profile-info}{หน้าโปรไฟล์ แท็บข้อมูล}
\appscreen{profile-history}{หน้าโปรไฟล์ แท็บประวัติ}
\appscreen{profile-events}{หน้าโปรไฟล์ แท็บกิจกรรมของฉัน}
\appscreen{profile-security}{หน้าโปรไฟล์ แท็บความปลอดภัย}

\section{หน้าจอของผู้ดูแลระบบ}

\subsection{การเข้าสู่ระบบและแดชบอร์ด}

ภาพที่ \ref{fig:app_admin-login} แสดงหน้าเข้าสู่ระบบของผู้ดูแลระบบ (UC-11) และภาพที่ \ref{fig:app_admin-dashboard} แสดงแดชบอร์ดที่สรุปสถิติของระบบ พร้อมเมนูด้านข้างที่แสดงจำนวนรายการรอตรวจสอบ (UC-14)

\appscreen[0.7]{admin-login}{หน้าเข้าสู่ระบบของผู้ดูแลระบบ}
\appscreen{admin-dashboard}{หน้าแดชบอร์ดผู้ดูแลระบบ}

\subsection{การจัดการสถานที่และกิจกรรม}

ภาพที่ \ref{fig:app_admin-places} ถึง \ref{fig:app_admin-events} แสดงรายการและแบบฟอร์มเพิ่มหรือแก้ไขข้อมูลสถานที่ และรายการกิจกรรม ซึ่งผู้ดูแลระบบเพิ่ม แก้ไข ลบ และเปิดหรือปิดการแสดงผลได้ (UC-06.1 ถึง UC-06.3 และ UC-09) หน้าจัดการสถานที่กรองรายการตามสถานะ อำเภอ และตำบลได้ แต่ละการ์ดแสดงวันที่ซิงก์ล่าสุด และแถบด้านบนมีปุ่มซิงก์จาก Google ตั้งค่าซิงก์อัตโนมัติ (UC-25) เพิ่มจาก Google Maps และนำเข้าตามพื้นที่ (UC-22)

\appscreen{admin-places}{หน้าจัดการสถานที่}
\appscreen{admin-place-form}{หน้าฟอร์มเพิ่มและแก้ไขข้อมูลสถานที่}
\appscreen{admin-events}{หน้าจัดการกิจกรรม}

\subsection{คิวตรวจสอบ}

ภาพที่ \ref{fig:app_admin-reports} และ \ref{fig:app_admin-event-requests} แสดงคิวรายงานข้อมูลคลาดเคลื่อนและคิวคำขอเพิ่มกิจกรรมจากผู้ใช้ที่ผู้ดูแลระบบตรวจสอบก่อนข้อมูลปรากฏในระบบ (UC-21)

\appscreen{admin-reports}{หน้ารายงานข้อมูล}
\appscreen{admin-event-requests}{หน้าคำขอกิจกรรม}

\subsection{QR ของรางวัล และการแลกที่เคาน์เตอร์}

ภาพที่ \ref{fig:app_admin-qr} และ \ref{fig:app_admin-rewards} แสดงการจัดการ QR Code และของรางวัล ส่วนภาพที่ \ref{fig:app_admin-redeem} แสดงหน้าที่เจ้าหน้าที่ใช้แลกของรางวัลให้ผู้ใช้ที่เคาน์เตอร์และยกเลิกรายการแลก (UC-06.5 และ UC-08)

\appscreen{admin-qr}{หน้าจัดการ QR Code}
\appscreen{admin-rewards}{หน้าจัดการของรางวัล}
\appscreen{admin-redeem}{หน้าแลกของรางวัลที่เคาน์เตอร์}

\subsection{ผู้ใช้และแบบสอบถาม}

ภาพที่ \ref{fig:app_admin-users} แสดงหน้าจัดการบัญชีผู้ใช้ (UC-23) และภาพที่ \ref{fig:app_admin-survey} แสดงหน้าสรุปผลแบบสอบถามความพึงพอใจ (UC-24)

\appscreen{admin-users}{หน้าจัดการผู้ใช้}
\appscreen{admin-survey}{หน้าผลความพึงพอใจ}

% =============================
% ภาคผนวก ข
% =============================
% ข้อความ Prompt อยู่ในไฟล์ backmatter/prompts/*.txt คัดจากโค้ด chatbot-service
% (ฉบับปัจจุบัน: คอมมิต 2333f89, ฉบับที่ใช้ทดลอง: คอมมิต c2b8db5 และ evals/trip_prompt/prompts.py คอมมิต 5aaa928)
% ถ้า Prompt ในโค้ดเปลี่ยน ให้คัดลอกใหม่และแก้เลขคอมมิตในย่อหน้าแรกของภาคผนวกนี้
\appendiceschapter{Prompt ที่ใช้ในระบบปัญญาประดิษฐ์}

ภาคผนวกนี้แสดง Prompt ของ 3 ส่วนที่บทที่ 4 อ้างถึง ได้แก่ ตัวสร้างแผนการเดินทาง กลไก LLM-as-a-Judge และแชตบอต คัดจากโค้ดของ AI Service ฉบับคอมมิต 2333f89 (22 ตุลาคม 2569) ยกเว้นฉบับที่ใช้ในการทดลองเปรียบเทียบ Prompt ซึ่งระบุคอมมิตไว้ในหัวข้อนั้น Prompt ทุกชุดเป็นแม่แบบที่ระบบเติมค่าตอนทำงาน ข้อความในวงเล็บปีกกา เช่น \{trip\_duration\_days\} คือตำแหน่งที่ระบบแทนด้วยค่าจริงจากเงื่อนไขของผู้ใช้หรือจากฐานข้อมูล ส่วนข้อความในวงเล็บที่เขียนเป็นภาษาอังกฤษท้ายบรรทัด เช่น (one line per candidate place) เป็นคำอธิบายที่ผู้วิจัยเพิ่มเพื่อการอ่าน ไม่ได้ส่งให้ LLM Prompt ของตัวสร้างแผนและ Judge เขียนเป็นภาษาอังกฤษ แต่สั่งให้เขียนข้อความที่แสดงแก่ผู้ใช้เป็นภาษาไทย ส่วน Prompt ของแชตบอตเขียนเป็นภาษาไทยทั้งหมด

\section{Prompt ของตัวสร้างแผนการเดินทาง}

ตัวสร้างแผนส่งข้อความ 2 ส่วนให้ LLM ส่วนแรกเป็นข้อความระบบ (System Message) ว่า ``You are a helpful travel assistant. Output JSON only.'' ส่วนที่สองเป็น Prompt ตามโปรแกรมที่ \ref{lst:trip_current} โดยตั้งค่า temperature เป็น 0.2 และกำหนดโครงสร้าง JSON ของผลลัพธ์ หัวข้อ [FEEDBACK] ท้าย Prompt จะปรากฏเฉพาะรอบที่สร้างแผนใหม่หลัง Judge หรือ Scheduler พบปัญหา

\subsection{ฉบับที่ใช้งานปัจจุบัน}

\lstinputlisting[caption={Prompt ของตัวสร้างแผนการเดินทาง ฉบับที่ใช้งานปัจจุบัน},label={lst:trip_current}]{backmatter/prompts/trip-current.txt}

ค่าที่ระบบเติมในโปรแกรมที่ \ref{lst:trip_current} ที่ต้องอธิบายเพิ่ม มีดังนี้ \{meals\} เป็นข้อความบอกมื้ออาหารของแต่ละวันตามช่วงเวลาเดินทาง เช่น ``each day has exactly lunch and dinner -- pick one "ร้านอาหาร" per meal ...'' และเพิ่มคำสั่งเลือกร้านอาหารเช้าใกล้ที่พักเมื่อผู้ใช้เลือก ``อาหารเช้า'' บรรทัด Last day ปรากฏเฉพาะเมื่อวันสุดท้ายไม่ได้จบที่ที่พัก โดยบอกจุดจบ เวลาที่ต้องถึง และให้เลือกสถานที่ที่อยู่ระหว่างทาง ในแต่ละบรรทัดของสถานที่ผู้สมัคร Type แสดงเมื่อสถานที่มีหมวดย่อย Price แสดงเฉพาะร้านอาหารและคาเฟ่ Visit คือเวลาแวะโดยประมาณตามประเภทสถานที่และจังหวะการท่องเที่ยว และ Hours แสดงเฉพาะเวลาทำการของวันในสัปดาห์ที่อยู่ในช่วงทริป ส่วน \{pace\_instruction\} และ \{budget\_instruction\} เลือกตามจังหวะการท่องเที่ยวและงบประมาณดังโปรแกรมที่ \ref{lst:trip_dynamic}

\lstinputlisting[caption={ข้อความของ \{pace\_instruction\} และ \{budget\_instruction\} ตามค่าที่ผู้ใช้เลือก},label={lst:trip_dynamic}]{backmatter/prompts/trip-dynamic.txt}

\subsection{ฉบับที่ใช้ในการทดลองเปรียบเทียบ Prompt}

การทดลองเปรียบเทียบ Prompt ในบทที่ 4 เริ่มรันเมื่อวันที่ 7 ตุลาคม 2569 ก่อนที่ Prompt จะถูกปรับเป็นฉบับปัจจุบัน ผู้วิจัยระบุฉบับที่ใช้จากค่าแฮช SHA-256 ของไฟล์ที่ชุดทดลองบันทึกไว้พร้อมผล ซึ่งตรงกับไฟล์ llm\_extractor.py ของคอมมิต c2b8db5 และไฟล์ prompts.py ของคอมมิต 5aaa928 Prompt ที่ใช้งานจริงขณะนั้น (v0\_baseline) แสดงในโปรแกรมที่ \ref{lst:trip_experiment} ซึ่งยังไม่มีประเภทอาหาร บรรทัดกำหนดมื้ออาหารของแต่ละวัน จุดจบทริป เวลาแวะ ระดับราคา และ rationale ขณะนั้น \{budget\_instruction\} ของงบหรูหราคือ ``- BUDGET (Luxury): Focus ONLY on providing the most premium, high-end, and exclusive experiences.'' ส่วนงบอื่นและ \{pace\_instruction\} เหมือนโปรแกรมที่ \ref{lst:trip_dynamic} (ในโปรแกรมที่ \ref{lst:trip_dynamic} คำว่า (coffee emoji) แทนอีโมจิรูปถ้วยกาแฟในโค้ดจริง เพราะฟอนต์ของโค้ดไม่มีอักขระนี้)

\lstinputlisting[caption={Prompt ของตัวสร้างแผนการเดินทาง ฉบับที่ใช้ในการทดลอง (v0\_baseline)},label={lst:trip_experiment}]{backmatter/prompts/trip-experiment.txt}

Prompt ทางเลือกทั้ง 3 แบบสร้างจาก v0\_baseline โดยแก้เพียงจุดเดียว และโค้ดของชุดทดลองตรวจว่าข้อความที่แก้พบเพียงครั้งเดียว หากไม่พบ การทดลองจะหยุดทันทีแทนที่จะทดสอบ Prompt ที่ต่างจากที่ตั้งใจ v1\_rules เพิ่มกฎข้อ 9--11 ในโปรแกรมที่ \ref{lst:trip_v1} ต่อท้ายหัวข้อ [STRICT RULES] v2\_context เปลี่ยนบรรทัดของสถานที่ผู้สมัครแต่ละแห่งเป็นรูปแบบในโปรแกรมที่ \ref{lst:trip_v2} ซึ่งเพิ่มจำนวนรีวิว แท็กไม่เกิน 5 แท็ก คำอธิบายย่อไม่เกิน 90 ตัวอักษร และเครื่องหมาย [MUST-GO] ของสถานที่บังคับ และ v3\_plan\_first เพิ่มขั้นตอนวางแผนก่อนเลือกสถานที่ในโปรแกรมที่ \ref{lst:trip_v3} ก่อนหัวข้อ [EXAMPLE JSON OUTPUT FORMAT] พร้อมเพิ่มคีย์ analysis ในตัวอย่าง JSON

\lstinputlisting[caption={{กฎที่ v1\_rules เพิ่มต่อท้ายหัวข้อ STRICT RULES}},label={lst:trip_v1}]{backmatter/prompts/trip-v1.txt}

\lstinputlisting[caption={รูปแบบบรรทัดของสถานที่ผู้สมัครใน v2\_context},label={lst:trip_v2}]{backmatter/prompts/trip-v2.txt}

\lstinputlisting[caption={ขั้นตอนวางแผนที่ v3\_plan\_first เพิ่มเข้าไป},label={lst:trip_v3}]{backmatter/prompts/trip-v3.txt}

\section{Prompt ของกลไก LLM-as-a-Judge}

Judge ส่งข้อความระบบว่า ``You are a meticulous travel-plan quality reviewer. Output JSON only.'' ร่วมกับ Prompt ในโปรแกรมที่ \ref{lst:judge} โดยตั้งค่า temperature เป็น 0 และกำหนดโครงสร้าง JSON ของผลลัพธ์ ข้อมูลในหัวข้อ [TRIP SETTINGS] และ [ITINERARY AS SCHEDULED] เช่น เวลาว่างของแต่ละวัน เวลากลับที่พัก และสถานที่ที่ตรงกับความสนใจ คำนวณด้วยโค้ดก่อนส่ง Judge จึงไม่ต้องอนุมานเองจากตารางเวลา

\lstinputlisting[caption={Prompt ของกลไก LLM-as-a-Judge},label={lst:judge}]{backmatter/prompts/judge.txt}

\section{Prompt ของแชตบอต}

แชตบอตเรียก LLM 2 ขั้น ขั้นแรกเขียนคำถามใหม่ให้เข้าใจได้ด้วยตัวเองเมื่อมีประวัติการสนทนา ตามโปรแกรมที่ \ref{lst:chat_rewrite} โดยตั้งค่า temperature เป็น 0 และระบบใช้เฉพาะข้อความภายในแท็ก <query> หากไม่พบแท็กจะใช้คำถามเดิม ขั้นที่สองสร้างคำตอบด้วย System Prompt ในโปรแกรมที่ \ref{lst:chat_system} ตามด้วยประวัติการสนทนาไม่เกิน 4 ข้อความ และคำถามของผู้ใช้ โดยตั้งค่า temperature เป็น 0.1 ทั้งนี้ ลำดับเลขข้อในหัวข้อ [กฎเหล็ก] เป็นไปตามโค้ดจริง ซึ่งข้อ 7 อยู่ท้ายสุด

\lstinputlisting[caption={Prompt สำหรับเขียนคำถามใหม่ (Query Rewriting)},label={lst:chat_rewrite}]{backmatter/prompts/chat-rewrite.txt}

\lstinputlisting[caption={System Prompt สำหรับสร้างคำตอบของแชตบอต},label={lst:chat_system}]{backmatter/prompts/chat-system.txt}

\{context\} ในโปรแกรมที่ \ref{lst:chat_system} คือข้อมูลที่ค้นคืนได้จากตาราง places, knowledge\_base และ events จัดรูปแบบตามโปรแกรมที่ \ref{lst:chat_context} รายการละหนึ่งชุดคั่นด้วย ``---'' รายละเอียดของแต่ละรายการตัดไว้ไม่เกิน 700 ตัวอักษร และหากค้นไม่พบรายการใดเลย ระบบใส่ข้อความ ``ไม่มีข้อมูลที่ตรงกับคำถามในฐานข้อมูล'' แทน

\lstinputlisting[caption={รูปแบบข้อมูลบริบทของสถานที่ ความรู้ทั่วไป และกิจกรรม},label={lst:chat_context}]{backmatter/prompts/chat-context.txt}
